% Original: PDF strana 8, donji slajd.
\slajdid{03-s016}
\begin{frame}[t,label=03-s016]{Koder prioriteta: potreba za signalom GS}
\setlength{\parskip}{0pt}
% element: 03-s016:e01
Pored multipleksera (posebne vrste kodera), u složenim digitalnim
i procesorskim sistemima često se koristi \textbf{koder prioriteta}. Koduje \textbf{najveći indeks aktivnog ulaza}.
Za četiri ulaza potrebna su dva izlazna bita, \(A_1,A_0\).
\par\vspace{2mm}
\begin{columns}[c,onlytextwidth]
\begin{column}{.44\textwidth}\centering
\Oznaka{Funkcionalna tabela}
\par\vspace{1.5mm}
% element: 03-s016:e02
{\fontsize{10}{12}\selectfont\begingroup
\setlength{\tabcolsep}{3pt}\renewcommand{\arraystretch}{1.12}
\par\noindent
\begin{tabularx}{\linewidth}{*{6}{>{\centering\arraybackslash}X}}
\toprule I3&I2&I1&I0&A1&A0\\\midrule
1&X&X&X&1&1\\\cmidrule(lr){1-6}
0&1&X&X&1&0\\\cmidrule(lr){1-6}
0&0&1&X&0&1\\\cmidrule(lr){1-6}
0&0&0&1&0&0\\\cmidrule(lr){1-6}
0&0&0&0&0&0\\\bottomrule
\end{tabularx}\par
\endgroup
}
\end{column}
\begin{column}{.52\textwidth}
% element: 03-s016:e03
\textbf{Dve poslednje vrste daju isti kod \(00\):}
\begin{itemize}
\item aktivan je samo ulaz \(I_0\), sa najnižim indeksom;
\item nijedan ulaz nije aktivan.
\end{itemize}
\par\vspace{1mm}
\Rezultat{Zato uvodimo dodatni izlaz \textcolor{DEisticanje}{\textbf{GS}}
(\emph{group select}): pokazuje da li postoji bar jedan aktivan ulaz,
odnosno da li ima šta da se koduje.}
\end{column}
\end{columns}
\end{frame}
